% ID: SEL-D0081-OA-E03
% Tanda S081.
Sean $P$ ser profesor de autoescuela y $M$ haber mejorado. Hay 60 profesores, de los cuales $60\cdot0{,}95=57$ mejoran y 3 no. Entre los otros 140 asistentes, $140\cdot0{,}80=112$ mejoran y 28 no.

\textbf{a)} No mejoran $3+28=31$ asistentes. Así, $P(M^c)=31/200=\boxed{0{,}155}$.

\textbf{b)} Condicionamos al grupo de $57+112=169$ personas que mejoran. De ellas, 112 no son profesores:
\[
 P(P^c\mid M)=\boxed{\frac{112}{169}}\approx0{,}6627.
\]
